\documentclass[11pt,a4paper]{article}
\usepackage[utf8]{inputenc}
\usepackage[T1]{fontenc}
\usepackage{lmodern}
\usepackage{hyperref}
\usepackage{graphicx}
\usepackage{geometry}
\geometry{margin=2.5cm}
\usepackage{xcolor}
\usepackage{tcolorbox}
\usepackage{sectsty}
\sectionfont{\color{blue!60!black}}

\title{\textbf{Architektura Systemu: CatFact Integration Service}}
\author{Jakub Filipiak}
\date{\today}

\begin{document}
\maketitle

\begin{abstract}
Dokumentacja techniczna dla uslugi integrujacej dane z zewnetrznego API (CatFact) z systemem lokalnym oraz emulowanym srodowiskiem chmurowym Microsoft Azure (Azurite). Projekt zaprojektowano z wykorzystaniem nowoczesnych wzorcow inzynierii oprogramowania w ekosystemie .NET.
\end{abstract}

\section{Zalozenia Architektoniczne}
Aplikacja zostala zaimplementowana w oparciu o szablon \textbf{.NET Worker Service}. Zapewnia to natywne wsparcie dla wstrzykiwania zaleznosci (Dependency Injection), zarzadzania cyklem zycia aplikacji oraz logowania.
Zewnetrzne API (\url{https://catfact.ninja/fact}) jest odpytywane ze stala czestotliwoscia 1 zapytania na sekunde przy uzyciu \texttt{PeriodicTimer}.

\section{Wzorce Projektowe i Siec}
\begin{itemize}
    \item \textbf{IHttpClientFactory}: Uzyty w celu unikniecia problemu wyczerpania gniazd sieciowych (socket exhaustion).
    \item \textbf{Polly (Resilience)}: Zaimplementowano polityke \textit{Retry} obslugujaca ewentualne problemy z dostepnoscia zewnetrznego endpointu.
    \item \textbf{Wydajne I/O}: Zapis do pliku \texttt{.txt} realizowany jest poprzez asynchroniczny, otwarty strumien (\texttt{StreamWriter}), redukujac narzut systemowy.
\end{itemize}

\section{Infrastruktura i Srodowisko Lokalne (Docker \& Azurite)}
Aby zapewnic pelna hermetycznosc srodowiska deweloperskiego, zastosowano konteneryzacje za pomoca \texttt{Docker} oraz \texttt{docker-compose}.
W ramach lokalnej chmury uruchamiany jest \textbf{Azurite} -- oficjalny emulator Microsoft Azure Storage.
Dzieki temu aplikacja moze bezpiecznie symulowac wysylanie kopii zapasowych pliku \texttt{.txt} do Azure Blob Storage bez koniecznosci podpinania rzeczywistej subskrypcji chmurowej przez osobe weryfikujaca zadanie.

\section{Infrastruktura jako Kod (IaC)}
W repozytorium znajduje sie wydzielony modul \texttt{/infrastructure} zawierajacy pliki \texttt{Terraform}. Sluza one wylacznie celom demonstracyjnym, obrazujac gotowosc projektu do wdrozenia w rzeczywistym srodowisku chmurowym (np. Azure Container Apps, Azure Storage Account).

\end{document}